\documentclass[11pt,a4paper]{article}
\usepackage[margin=25mm]{geometry}
\usepackage{iftex}
\ifPDFTeX
  \usepackage[T2A]{fontenc}
  \usepackage[utf8]{inputenc}
\else
  \usepackage{fontspec}
  \setmainfont{cmunrm.otf}[
    BoldFont=cmunbx.otf,
    ItalicFont=cmunti.otf,
    BoldItalicFont=cmunbi.otf,
    Ligatures=TeX]
\fi
\usepackage[english,russian]{babel}
\usepackage{amsmath,amssymb,amsthm}
\usepackage{mathrsfs}
\usepackage{microtype}
\usepackage{enumitem}
\usepackage[unicode]{hyperref}
\hypersetup{
  pdftitle={Неограниченность индексов орбитальных рядов и редуктивности конечных квандлов},
  pdfauthor={Danil Pokulevskii},
  pdflang={ru-RU},
  colorlinks=true,linkcolor=blue,citecolor=blue,urlcolor=blue}
\setlength{\emergencystretch}{2em}
\setlist[enumerate]{itemsep=2pt,topsep=5pt}
\frenchspacing
\newtheorem{theorem}{Теорема}[section]
\newtheorem{lemma}[theorem]{Лемма}
\newtheorem{proposition}[theorem]{Предложение}
\newtheorem{corollary}[theorem]{Следствие}
\theoremstyle{definition}
\newtheorem{definition}[theorem]{Определение}
\theoremstyle{remark}
\newtheorem{remark}[theorem]{Замечание}
\renewcommand{\proofname}{Доказательство}
\newcommand{\R}{\mathcal R}
\newcommand{\LR}{\mathcal{LR}}
\newcommand{\tOS}{t\mathcal{OS}}
\newcommand{\Inn}{\operatorname{Inn}}
\newcommand{\Orb}{\operatorname{Orb}}
\newcommand{\cl}{\operatorname{class}}
\newcommand{\ad}{\operatorname{ad}}
\newcommand{\Z}{\mathbb Z}
\newcommand{\F}{\mathbb F}
\newcommand{\conj}{\triangleright}

\title{Неограниченность индексов орбитальных рядов\\
и редуктивности конечных квандлов}
\author{Danil Pokulevskii\\[3pt]
  \small Новосибирский государственный технический университет\\
  \small Новосибирск, Россия\\
  \small\href{mailto:555tery@gmail.com}{\texttt{555tery@gmail.com}}}
\date{7 октября 2026 года\\[2pt]\small Рабочая версия препринта}

\begin{document}
\maketitle
\begin{abstract}
Для конечного квандла обозначим через \(k\), \(j\), \(i\)
минимальные положительные индексы локальной редуктивности,
тривиализации орбитальных рядов и редуктивности соответственно.
Показывается, что ни \(j\) нельзя ограничить функцией от \(k\),
ни \(i\) нельзя ограничить функцией от \(j\).
Для каждого \(n\ge3\) построен конечный квандл \(Q_n\) с
\(k(Q_n)=2\) и \(j(Q_n)\ge\lfloor\log_3n\rfloor+1\).
Для каждого \(m\ge2\) построен конечный квандл \(P_m\) с
\((k(P_m),j(P_m),i(P_m))=(2,2,m)\).
Первое семейство получается из группы Ноче--Трейси--Траустасона
посредством фильтрации по размеру метки; второе --- из конечных
факторов редуцированной свободной группы, сохраняющих точный
класс нильпотентности. Все переходы от групп к квандлам
и сдвиг между классом внутренней группы и индексом редуктивности
доказываются явно. Тем самым получены отрицательные ответы
на обе части вопроса~5.25 Бонатто--Кранс--Насыбуллова--Уитни.
\end{abstract}

\noindent\textbf{Ключевые слова:} конечный квандл, локальная редуктивность,
орбитальный ряд, редуцированная свободная группа, элемент Энгеля.

\section{Введение и основные результаты}

В работе~\cite{BCNW} введены классы \(\LR_n\), \(\tOS_n\)
и изучена их связь с классами \(\R_n\).
Для конечных квандлов эти условия конечной глубины эквивалентны
на уровне объединений классов \cite[следствие~5.19]{BCNW}.
Вопрос~5.25 спрашивает, можно ли ограничить глубину орбитальных
рядов через индекс локальной редуктивности, а индекс редуктивности ---
через глубину орбитальных рядов.
Нумерация результатов~\cite{BCNW} здесь относится к
arXiv v3 от 18 августа 2020 года.

Для квандла, у которого соответствующие индексы существуют,
полагаем
\[
 \begin{aligned}[c]
 k(Q)&=\min\{n\ge1:Q\in\LR_n\},\\
 j(Q)&=\min\{n\ge1:Q\in\tOS_n\},\\
 i(Q)&=\min\{n\ge1:Q\in\R_n\}.
 \end{aligned}
\]
Определения классов приведены в разделе~\ref{sec:definitions}.

\begin{theorem}\label{thm:first}
Для каждого целого \(n\ge3\) существует конечный квандл \(Q_n\),
для которого
\[
 k(Q_n)=2,\qquad j(Q_n)\ge\lfloor\log_3n\rfloor+1.
\]
\end{theorem}

\begin{theorem}\label{thm:second}
Для каждого целого \(m\ge2\) существует конечный квандл \(P_m\),
для которого
\[
 (k(P_m),j(P_m),i(P_m))=(2,2,m).
\]
Квандл \(P_m\) можно выбрать с ровно \(m\) внутренними орбитами,
каждая из которых тривиальна как квандл.
\end{theorem}

\begin{corollary}\label{cor:nobounds}
Не существует функции \(f:\mathbb N_{>0}\to\mathbb N_{>0}\),
для которой \(j(Q)\le f(k(Q))\) для всех конечных квандлов
с определёнными индексами. Также не существует такой функции \(g\),
для которой \(i(Q)\le g(j(Q))\) для всех этих квандлов.
\end{corollary}
\begin{proof}
Первое семейство имеет \(k=2\) и неограниченные \(j\).
Второе имеет \(j=2\) и неограниченные \(i\).
\end{proof}

Доказательства теорем~\ref{thm:first} и~\ref{thm:second}
даны в разделах~\ref{sec:first} и~\ref{sec:second}.
Используются две известные групповые конструкции:
группа Ноче--Трейси--Траустасона~\cite{NTT}
и редуцированная свободная группа в изложении Дарне~\cite{Darne}.
Групповые входные результаты указаны отдельно от выводимых
здесь утверждений о квандлах.

\section{Определения и критерий редуктивности}\label{sec:definitions}

Квандл понимается в левом соглашении~\cite{BCNW}:
это непустое множество \(Q\) с операцией \(\conj\), для которой
\[
 a\conj a=a,\qquad L_a:b\mapsto a\conj b\ \text{биективно},
 \qquad a\conj(b\conj c)=(a\conj b)\conj(a\conj c).
\]
Группа \(\Inn(Q)=\langle L_a:a\in Q\rangle\)
называется внутренней группой.
Под-квандлы замкнуты также относительно обратных левых трансляций.
Тривиальная операция задаётся равенством \(a\conj b=b\).

Обозначим через
\[
 T_n(a;c_1,\ldots,c_n)
 =((a\conj c_1)\conj\cdots)\conj c_n
\]
левоассоциированное выражение с \(n\) операциями.
Класс \(\R_n\) задаётся тождеством
\[
 T_n(a;c_1,\ldots,c_n)=T_n(b;c_1,\ldots,c_n)
\]
для всех аргументов; класс \(\LR_n\) --- тождеством
\[
 T_n(a;\underbrace{b,\ldots,b}_{n})=b.
\]

Орбитальный ряд строится рекурсивно:
\[
 Q_0=Q,\qquad Q_{r+1}=\Orb(x_r,Q_r),\qquad x_r\in Q_r.
\]
На каждом шаге используется внутренняя группа текущего под-квандла.
Класс \(\tOS_n\) состоит из квандлов, для которых любой такой
ряд достигает одноэлементного квандла не позднее шага \(n\).
Таким образом, \(j\) измеряет достижение одноэлементного квандла,
а не только стабилизацию ряда.
Все три класса с индексом \(1\) совпадают с классом тривиальных квандлов.
Кроме того, \(Q\in\tOS_2\) тогда и только тогда, когда
каждая внутренняя орбита \(Q\) тривиальна как квандл;
сама орбита при этом может содержать много элементов.

Для групп используем соглашения
\[
 [a,b]=a^{-1}b^{-1}ab,\qquad c^g=g^{-1}cg,\qquad
 \gamma_1(H)=H,\quad \gamma_{r+1}(H)=[\gamma_r(H),H].
\]
Класс тривиальной группы равен нулю.
В остальных случаях точный класс \(c\) означает
\(\gamma_{c+1}(H)=1\), \(\gamma_c(H)\ne1\).

\begin{lemma}\label{lem:orbits}
Если \(S\) --- непустое подмножество группы, замкнутое относительно
сопряжений своими элементами и их обратными, то операция
\(a\conj b=a^{-1}ba\) делает \(S\) квандлом. Для \(a\in S\)
\[
 \Orb(a,S)=a^{\langle S\rangle}.
\]
\end{lemma}
\begin{proof}
Идемпотентность и биективность непосредственны, а
\[
 (a^{-1}ba)^{-1}(a^{-1}ca)(a^{-1}ba)=a^{-1}b^{-1}cba
\]
доказывает самодистрибутивность.
Левые трансляции и их обратные суть сопряжения элементами
\(S\) и их обратными. Они порождают действие \(\langle S\rangle\),
откуда формула для орбиты.
\end{proof}

\begin{proposition}\label{prop:criterion}
Для любого непустого квандла \(Q\) и \(n\ge1\)
\[
 Q\in\R_n\quad\Longleftrightarrow\quad
 \gamma_n(\Inn(Q))=1.
\]
В частности, если \(\Inn(Q)\) имеет точный конечный класс \(c\),
то \(i(Q)=c+1\).
\end{proposition}
\begin{proof}
Положим \(\mathcal L(Q)=\{L_a:a\in Q\}\).
Из самодистрибутивности следует
\[
 L_{a\conj b}=L_aL_bL_a^{-1}.
\]
Поэтому \(a\mapsto L_a\) есть эпиморфизм квандлов,
если на \(\mathcal L(Q)\) использовать операцию \(s\conj t=sts^{-1}\).
Это направление сопряжения противоположно соглашению
леммы~\ref{lem:orbits}.

При \(n\ge1\) тождество \(\R_{n+1}\) для \(Q\)
эквивалентно тождеству \(\R_n\) для \(\mathcal L(Q)\):
равенство двух выражений после последней операции
для всех последних аргументов означает равенство их левых трансляций.
Также \(Q\in\R_1\) тогда и только тогда, когда
\(\mathcal L(Q)\) одноэлементен.
Действительно, если все \(L_a\) равны \(f\), то
идемпотентность даёт \(f(a)=a\) для всех \(a\).
Для итераций \(\mathcal L_0(Q)=Q\),
\(\mathcal L_{r+1}(Q)=\mathcal L(\mathcal L_r(Q))\) получаем
\[
 Q\in\R_n\quad\Longleftrightarrow\quad
 |\mathcal L_n(Q)|=1.
\]

Пусть \(H=\Inn(Q)\). Множество \(\mathcal L(Q)\) порождает \(H\)
и сохраняется сопряжениями \(H\), поскольку
\(hL_ah^{-1}=L_{h(a)}\).
Действие \(H\) на этом множестве имеет образ
\(\Inn(\mathcal L(Q))\) и ядро
\(C_H(\mathcal L(Q))=Z(H)\).
Следовательно,
\[
 \Inn(\mathcal L(Q))\cong H/Z(H).
\]
Для верхнего центрального ряда \(Z_0(H)=1\),
\(Z_{r+1}(H)/Z_r(H)=Z(H/Z_r(H))\)
индукция даёт
\[
 \Inn(\mathcal L_r(Q))\cong H/Z_r(H).
\]
Квандл \(\mathcal L_{n-1}(Q)\) тривиален тогда и только тогда,
когда \(\mathcal L_n(Q)\) одноэлементен; поэтому
последнее равенство эквивалентно \(H=Z_{n-1}(H)\).

Наконец, \(H=Z_c(H)\) эквивалентно \(\gamma_{c+1}(H)=1\).
Это следует индукцией по \(c\), начиная с \(c=0\):
\(\gamma_{c+1}(H)=1\) эквивалентно
\(\gamma_c(H)\le Z(H)\), то есть
\(\gamma_c(H/Z(H))=1\).
Применение индукции к \(H/Z(H)\) даёт \(H=Z_c(H)\).
Подставляя \(c=n-1\), получаем критерий, включая \(n=1\).
\end{proof}

\begin{remark}
Критерий соответствует \cite[предложение~3.2]{BCNW},
но записан через исчезновение члена ряда, чтобы отличить
верхнюю оценку класса от точного класса.
\end{remark}

\section{Первое семейство: неограниченная глубина}
\label{sec:first}

\subsection{Групповые входные данные}

Используем группу \(\mathscr G\) над \(\F_2\), построенную
в~\cite{NTT}. Её генераторы обозначим
\[
 X=1+\ad(x),\quad U_A=1+\ad(u_A),\quad
 V_A=1+\ad(v_A),\quad W_A=1+\ad(w_A),
\]
где \(A\) --- непустое конечное подмножество \(\mathbb N\).
Операторы и их действие определены в~\cite[раздел~2]{NTT}.
Все указанные генераторы --- инволюции.
Каждое из трёх семейств \(U,V,W\) порождает элементарно абелеву группу.
Если \(A\cap B=\varnothing\), обозначим \(A\sqcup B=A\cup B\);
иначе положим \(A\sqcup B=\varnothing\).
Пустая метка соответствует единице.
По \cite[лемма~2.3]{NTT}
\begin{equation}\label{eq:nttcomm}
 \begin{aligned}[c]
 [U_A,V_B]&=U_{A\sqcup B},&
 [V_A,W_B]&=W_{A\sqcup B},\\
 [W_A,U_B]&=V_{A\sqcup B},&
 [W_A,X]&=U_A,
 \end{aligned}
\end{equation}
а \(X\) коммутирует со всеми \(U_A,V_A\).
По \cite[предложение~2.4]{NTT} каждый элемент имеет
единственную нормальную форму \(X^\epsilon uvw\),
где \(\epsilon\in\{0,1\}\), а \(u,v,w\) принадлежат
соответствующим элементарно абелевым группам.
Наконец, в доказательстве \cite[теорема~2.5]{NTT} установлено
\begin{equation}\label{eq:pairengel}
 [[X^g,X],X]=1\qquad(g\in\mathscr G).
\end{equation}

Каждый \(U_A\) нетривиален. Чтобы это увидеть в исходной
операторной конструкции, выбираем одноэлементное \(B\),
не пересекающееся с \(A\): оператор \(\ad(u_A)\)
переводит \(w_B\) в ненулевой \(v_{A\cup B}\).

Для \(n\ge1\) положим
\[
 G_n=\langle X,U_A,V_A,W_A:\varnothing\ne A\subseteq[n]\rangle
 \le\mathscr G,\qquad [n]=\{1,\ldots,n\}.
\]
При сборке слова в нормальную форму новые метки получаются
только как \(\sqcup\), поэтому остаются подмножествами \([n]\).
Отсюда \(G_n\) конечна и
\begin{equation}\label{eq:sizefirst}
 |G_n|\le2^{\,1+3(2^n-1)}.
\end{equation}
Здесь ограничиваются сами метки, а не только их мощности.

\subsection{Фильтрация и нормальные замыкания}

Для целого \(s\ge1\) введём подгруппу
\[
 F_s=\langle X,\ U_A\ (|A|\ge s),\
 V_A\ (|A|\ge2s),\ W_A\ (|A|\ge3s)\rangle\le G_n,
\]
где используются только метки в \([n]\).
Её элементы имеют в точности нормальные формы с
\[
 u\in\langle U_A:|A|\ge s\rangle,\quad
 v\in\langle V_A:|A|\ge2s\rangle,\quad
 w\in\langle W_A:|A|\ge3s\rangle.
\]
Действительно, единственные поправки при сборке удовлетворяют
\[
 \begin{aligned}[c]
 [X,W_{\ge3s}]&\subseteq U_{\ge3s},&
 [U_{\ge s},V_{\ge2s}]&\subseteq U_{\ge3s},\\
 [W_{\ge3s},U_{\ge s}]&\subseteq V_{\ge4s},&
 [V_{\ge2s},W_{\ge3s}]&\subseteq W_{\ge5s}.
 \end{aligned}
\]
Порог не понижается: непересекающиеся метки дают сумму размеров,
а пересекающиеся --- единицу. Обозначения \(U_{\ge t}\)
и аналогичные далее означают подгруппы, порождённые метками
соответствующего размера.

\begin{lemma}\label{lem:filtration}
Для каждого \(s\ge1\)
\[
 \langle X^{F_s}\rangle=F_{3s}.
\]
Кроме того, \(\langle X^{G_n}\rangle=F_1\).
\end{lemma}
\begin{proof}
Сначала проверим \(F_{3s}\trianglelefteq F_s\).
В таблице строки соответствуют генераторам \(F_{3s}\),
столбцы --- генераторам \(F_s\); запись указывает слой,
содержащий возможный ненулевой коммутатор:
\[
 \begin{array}{c|cccc}
 &X&U_{\ge s}&V_{\ge2s}&W_{\ge3s}\\ \hline
 X&1&1&1&U_{\ge3s}\\
 U_{\ge3s}&1&1&U_{\ge5s}&V_{\ge6s}\\
 V_{\ge6s}&1&U_{\ge7s}&1&W_{\ge9s}\\
 W_{\ge9s}&U_{\ge9s}&V_{\ge10s}&W_{\ge11s}&1
 \end{array}
\]
Все записи входят в \(F_{3s}\).
Так как генераторы --- инволюции, проверка обеспечивает
инвариантность при сопряжении ими в обоих направлениях.
Подгруппа содержит \(X\), поэтому
\(\langle X^{F_s}\rangle\le F_{3s}\).

Для обратного включения положим \(K=\langle X^{F_s}\rangle\).
Она нормальна в \(F_s\) и содержит \(X\).
При \(|A|\ge3s\) имеем \(W_A\in F_s\), откуда
\([W_A,X]=U_A\in K\).
Каждую метку \(T\) размера не меньше \(6s\)
можно разбить на \(A\sqcup B\) с \(|A|,|B|\ge3s\);
тогда \([W_A,U_B]=V_T\in K\).
При \(|T|\ge9s\) берём разбиение
\(|A|\ge6s,\ |B|\ge3s\).
Предыдущий шаг даёт \(V_A\in K\), поэтому получаем
\([V_A,W_B]=W_T\in K\).
Это все генераторы \(F_{3s}\).

Для начального шага проверяем \(F_1\trianglelefteq G_n\).
Возможные ненулевые коммутаторы по образующим лежат в слоях
\[
 \begin{aligned}[c]
 [X,W_{\ge1}]&\subseteq U_{\ge1},&
 [U_{\ge1},V_{\ge1}]&\subseteq U_{\ge2},\\
 [U_{\ge1},W_{\ge1}]&\subseteq V_{\ge2},&
 [V_{\ge2},W_{\ge1}]&\subseteq W_{\ge3},\\
 [W_{\ge3},U_{\ge1}]&\subseteq V_{\ge4},&
 [W_{\ge3},V_{\ge1}]&\subseteq W_{\ge4},\\
 [W_{\ge3},X]&\subseteq U_{\ge3}.
 \end{aligned}
\]
Все они входят в \(F_1\), остальные равны единице
или получаются обращением коммутатора.
Значит, \(\langle X^{G_n}\rangle\le F_1\).
Обратно, из \([W_A,X]=U_A\) нормальное замыкание содержит
все \(U_A\); из разбиений на две непустые части и
\([W_A,U_B]\) --- все \(V_T\) с \(|T|\ge2\);
из разбиений \(|A|\ge2,\ |B|\ge1\) и
\([V_A,W_B]\) --- все \(W_T\) с \(|T|\ge3\).
\end{proof}

\subsection{Квандл и его глубина}

\begin{proof}[Доказательство теоремы~\ref{thm:first}]
Пусть \(C_n=X^{G_n}\), и на \(Q_n=C_n\)
зададим операцию \(a\conj b=a^{-1}ba\).
Это конечный квандл по лемме~\ref{lem:orbits}.
Для \(a=X^g,\ b=X^h\) сопряжение на \(h^{-1}\)
переводит пару в \(X^{gh^{-1}},X\).
Поэтому из \eqref{eq:pairengel} получаем
\([[a,b],b]=1\) для всех \(a,b\in C_n\).
Перестановка \(a,b\) даёт \([[b,a],a]=1\).
Так как \([b,a]=[a,b]^{-1}\), коммутатор \([a,b]\)
централизует обоих генераторов. Подгруппа \(\langle a,b\rangle\)
имеет класс не выше двух.
В частности, \(a^{-1}ba\) коммутирует с \(b\), откуда
\[
 (a\conj b)\conj b=b.
\]
Значит, \(Q_n\in\LR_2\).

При \(n\ge3\) выберем \(|A|=3\).
Тогда \(W_A\in F_1=\langle C_n\rangle\)
и \([W_A,X]=U_A\ne1\).
По лемме~\ref{lem:orbits} орбита \(X\) в \(Q_n\)
нетривиальна, так что \(Q_n\) не тривиален и \(k(Q_n)=2\).
Индекс \(j(Q_n)\) существует:
по \cite[следствие~5.19]{BCNW} конечный квандл из \(\LR_2\)
лежит в \(\tOS\), то есть каждый его орбитальный ряд
достигает одноэлементного квандла.
В конечном квандле длина такого ряда ограничена:
до одноэлементного члена все включения строги,
поскольку равенство соседних членов означало бы
связность текущего члена и невозможность дальнейшего уменьшения.
Поэтому \(Q_n\in\bigcup_{t\ge1}\tOS_t\).

Выбирая на каждом шаге якорь \(X\), положим
\[
 S_0=C_n,\qquad N_r=\langle S_r\rangle,\qquad
 S_{r+1}=X^{N_r}.
\]
Это орбитальный ряд по лемме~\ref{lem:orbits}.
Лемма~\ref{lem:filtration} даёт
\[
 N_r=F_{3^r}\quad(r\ge0),\qquad
 S_r=X^{F_{3^{r-1}}}\quad(r\ge1).
\]
Если \(r\ge1\) и \(3^r\le n\), берём \(|A|=3^r\).
Тогда \(W_A\in F_{3^{r-1}}\), поэтому
\(X,X^{W_A}\in S_r\), причём эти элементы различны,
так как \([W_A,X]=U_A\ne1\).
Следовательно, \(S_r\) не одноэлементен.
При \(r=\lfloor\log_3n\rfloor\) ряд не достигает
одноэлементного члена на шагах \(0,\ldots,r\),
и \(j(Q_n)\ge r+1\).
\end{proof}

\begin{remark}
В теореме~\ref{thm:first} установлена нижняя оценка.
Точные значения \(j(Q_n)\) и минимальные порядки
квандлов с заданными индексами здесь не устанавливаются.
\end{remark}

\section{Второе семейство: неограниченная редуктивность}
\label{sec:second}

\subsection{Редуцированная свободная группа и конечный фактор}

Пусть \(F(x_1,\ldots,x_m)\) --- свободная группа.
Определим редуцированную свободную группу
\[
 RF_m=F(x_1,\ldots,x_m)/
 \langle\!\langle [x_r,x_r^w]:
 1\le r\le m,\ w\in F(x_1,\ldots,x_m)\rangle\!\rangle.
\]
Каждый \(x_r\) коммутирует со всеми своими сопряжёнными.
По \cite[предложения~1.3 и~1.18]{Darne}
\[
 \gamma_{m+1}(RF_m)=1,\qquad
 \gamma_m(RF_m)=Z(RF_m)\cong\Z^{(m-1)!}.
\]
Здесь и далее нумерация~\cite{Darne} относится к arXiv v2.
Следовательно, \(RF_m\) имеет точный класс \(m\).
Выберем \(1\ne z_m\in\gamma_m(RF_m)\).

Обоснуем существование конечного фактора,
в котором \(z_m\) не исчезает.
Рассмотрим кольцо
\[
 \mathscr A_m=\Z\langle X_1,\ldots,X_m\rangle/I,
\]
где \(I\) порождён словами с повторением хотя бы одной буквы.
Аддитивный базис состоит из пустого слова и всех слов без повторений.
Для идеала \(J=(X_1,\ldots,X_m)\) имеем \(J^{m+1}=0\);
поэтому \(1+J\) --- группа с обратными
\[
 (1+u)^{-1}=\sum_{t=0}^{m}(-u)^t.
\]
По \cite[следствие~1.13]{Darne} отображение Магнуса
\[
 \mu:RF_m\hookrightarrow1+J,\qquad x_r\mapsto1+X_r
\]
инъективно.
В разложении \(\mu(z_m)-1=\sum_w c_ww\)
существует ненулевой целый коэффициент \(c_w\).
Выберем простое \(p\), не делящее этот коэффициент,
и редуцируем кольцо по модулю \(p\).
Группа \(1+\overline J\) конечна, поскольку \(\overline J\)
имеет конечный базис над \(\F_p\); её порядок равен \(p^{d_m}\), где
\[
 d_m=\sum_{t=1}^{m}\frac{m!}{(m-t)!}.
\]
Пусть \(A_m\) --- образ \(RF_m\) в этой группе,
а \(\phi_m:RF_m\twoheadrightarrow A_m\) --- эпиморфизм на образ.
Тогда \(\phi_m(z_m)\ne1\).
Эпиморфизмы сохраняют нижний центральный ряд:
\(\phi_m(\gamma_t(RF_m))=\gamma_t(A_m)\),
что следует индукцией из сохранения коммутаторов.
Значит, \(A_m\) имеет точный класс \(m\).

Чтобы классы выбранных генераторов были различны,
положим
\[
 \alpha:RF_m\to(\Z/2\Z)^m,\quad x_r\mapsto e_r,\qquad
 G_m=\operatorname{im}(\phi_m,\alpha)
 \le A_m\times(\Z/2\Z)^m.
\]
Отображение \(\alpha\) корректно, поскольку все отношения
лежат в коммутаторной подгруппе свободной группы.
Обозначим эпиморфизм на \(G_m\) через \(\psi_m\),
а образы \(x_r\) через \(a_r\).
Группа \(G_m\) конечна, порождается \(a_1,\ldots,a_m\)
и имеет точный класс \(m\):
она является фактором \(RF_m\) и сохраняет \(z_m\)
в первой координате.
Абелевы координаты \(e_r\) различны, поэтому
\(a_r,a_s\) не сопряжены при \(r\ne s\).

\subsection{Орбиты и три точных индекса}

\begin{proof}[Доказательство теоремы~\ref{thm:second}]
Фиксируем \(m\ge2\) и сокращённо пишем \(G=G_m\).
Положим
\[
 C_r=a_r^G,\qquad P_m=\bigcup_{r=1}^{m}C_r,\qquad
 u\conj v=u^{-1}vu.
\]
По лемме~\ref{lem:orbits} это конечный квандл.
Классы \(C_r\) попарно не пересекаются.
Отношения редуцированной свободной группы дают
\([a_r,a_r^g]=1\) для любого \(g\in G\).
Любые два элемента \(a_r^g,a_r^h\) коммутируют:
сопряжение на \(g^{-1}\) переводит их в
\(a_r,a_r^{hg^{-1}}\).
Таким образом, каждый \(C_r\) тривиален как квандл.

Рассмотрим гомоморфизм
\[
 \rho:G\to\operatorname{Sym}(P_m),\qquad
 \rho(g)(v)=gvg^{-1}.
\]
При обычной композиции перестановок это именно гомоморфизм;
отображения \(v\mapsto g^{-1}vg\) образуют антигомоморфизм.
Имеем \(L_u=\rho(u^{-1})\).
Поскольку \(P_m\) содержит все \(a_r\), оно порождает \(G\);
следовательно, \(\Inn(P_m)=\rho(G)\).
Его орбиты --- ровно \(C_1,\ldots,C_m\).
Любой орбитальный ряд достигает одного из \(C_r\)
на первом шаге и одноэлементного квандла на втором.
Значит, \(P_m\in\tOS_2\).
Квандл не тривиален: иначе все \(a_r\) попарно коммутировали бы,
а порождённая ими группа \(G\) была бы абелевой,
вопреки её классу \(m\ge2\).
Поэтому \(j(P_m)=2\).

Элемент \(u\conj v\) лежит в том же \(C_r\), что и \(v\),
поэтому коммутирует с \(v\). Отсюда
\((u\conj v)\conj v=v\), то есть \(P_m\in\LR_2\).
Нетривиальность даёт \(k(P_m)=2\).

Ядро \(\rho\) равно \(C_G(P_m)=Z(G)\),
поскольку \(P_m\) порождает \(G\).
Следовательно,
\[
 \Inn(P_m)\cong G/Z(G).
\]
Фактор \(G/Z(G)\) имеет точный класс \(m-1\).
Действительно, \(\gamma_m(G)\le Z(G)\), так что
\(\gamma_m(G/Z(G))=1\).
Если бы \(\gamma_{m-1}(G/Z(G))=1\), то
\(\gamma_{m-1}(G)\le Z(G)\), откуда \(\gamma_m(G)=1\),
противоречие.
По предложению~\ref{prop:criterion} получаем \(i(P_m)=m\).
Все требуемые равенства доказаны.
\end{proof}

\begin{remark}
Абелев фактор \((\Z/2\Z)^m\) нужен для утверждения о ровно
\(m\) различных орбитах. Для неограниченности \(i\) при \(j=2\)
достаточно самого \(A_m\).
Выбор простого \(p\) сделан после выбора \(z_m\);
сохранение произвольного \(z_m\) при любом заранее заданном \(p\)
не утверждается.
\end{remark}

\subsection{\texorpdfstring{Пример при \(m=2\)}{Пример при m=2}}

Для проверки сдвига индексов можно взять
\(G=UT_3(\F_2)\) с генераторами
\[
 a=1+E_{12},\qquad b=1+E_{23},\qquad z=1+E_{13}.
\]
Группа имеет порядок \(8\), класс \(2\) и центр \(\{1,z\}\).
Классы генераторов суть \(\{a,az\}\) и \(\{b,bz\}\).
На \(P=\{a,az,b,bz\}\) операция имеет таблицу
\[
 \begin{array}{c|cccc}
 \conj&a&az&b&bz\\ \hline
 a&a&az&bz&b\\
 az&a&az&bz&b\\
 b&az&a&b&bz\\
 bz&az&a&b&bz
 \end{array}
\]
Видны две тривиальные орбиты размера \(2\), а
\(\Inn(P)\cong G/Z(G)\) нетривиальна и абелева.
Следовательно, \((k(P),j(P),i(P))=(2,2,2)\).

\section{Область результата}

Оба препятствия возникают уже при минимальном нетривиальном
значении контролирующего индекса, равном \(2\).
Первое семейство состоит из одного группового класса сопряжения;
второе --- из объединения классов выбранных генераторов.
Это не полный квандл сопряжения на всей группе.
Поэтому \cite[следствие~5.24]{BCNW}, относящееся
к полному квандлу \(\operatorname{Conj}(G)\), не ограничивает
класс группы в конструкции \(P_m\).

Конечный фактор во втором семействе выбран так,
чтобы сохранить ненулевой коммутатор веса \(m\).
Произвольный конечный фактор может уменьшить класс
и не обеспечивает требуемого индекса.
Доказательства устанавливают существование семейств;
оптимальные порядки и точная глубина первого семейства
остаются за рамками статьи.

\begin{thebibliography}{9}
\bibitem{BCNW}
M.~Bonatto, A.~S.~Crans, T.~Nasybullov, G.~T.~Whitney.
\newblock Quandles with orbit series conditions.
\newblock \emph{Journal of Algebra}, \textbf{567} (2021), 284--309.
\newblock DOI:
\href{https://doi.org/10.1016/j.jalgebra.2020.09.026}{10.1016/j.jalgebra.2020.09.026}.
\newblock Нумерация:
\href{https://arxiv.org/abs/2004.10227v3}{arXiv:2004.10227v3}.

\bibitem{NTT}
M.~Noce, G.~Tracey, G.~Traustason.
\newblock A left \(3\)-Engel element whose normal closure is not nilpotent.
\newblock \emph{Journal of Pure and Applied Algebra},
\textbf{224}, no.~3 (2020), 1092--1101.
\newblock DOI:
\href{https://doi.org/10.1016/j.jpaa.2019.07.004}{10.1016/j.jpaa.2019.07.004}.
\newblock \href{https://arxiv.org/abs/1811.12074v1}{arXiv:1811.12074v1}.

\bibitem{Darne}
J.~Darn\'e.
\newblock Milnor invariants of braids and welded braids up to homotopy.
\newblock \emph{Algebraic \& Geometric Topology},
\textbf{24}, no.~3 (2024), 1277--1320.
\newblock DOI:
\href{https://doi.org/10.2140/agt.2024.24.1277}{10.2140/agt.2024.24.1277}.
\newblock Нумерация:
\href{https://arxiv.org/abs/1904.10677v2}{arXiv:1904.10677v2}.
\end{thebibliography}
\end{document}
